\documentclass[11pt,a4paper]{article}

\usepackage[margin=2.5cm]{geometry}
\usepackage[T1]{fontenc}
\usepackage{lmodern}
\usepackage{microtype}
\usepackage{amsmath,amssymb,mathtools}
\usepackage{booktabs}
\usepackage{xcolor}
\usepackage[most]{tcolorbox}
\usepackage{tikz}
\usetikzlibrary{arrows.meta,positioning,fit,backgrounds}
\PassOptionsToPackage{hyphens}{url}
\usepackage[colorlinks=true,linkcolor=blue!60!black,urlcolor=blue!60!black,citecolor=blue!60!black]{hyperref}
\usepackage[capitalise,noabbrev]{cleveref}

% ---------------------------------------------------------------------------
% Links to the two papers. Every equation reference below is clickable.
% ---------------------------------------------------------------------------
\newcommand{\urlPone}{https://arxiv.org/pdf/2511.03604}
\newcommand{\urlPtwo}{https://arxiv.org/pdf/2410.08263}
\newcommand{\Pone}{\href{\urlPone}{Buscicchio, Pozzoli et al.~(2025)}}
\newcommand{\Ptwo}{\href{\urlPtwo}{Buscicchio, Klein et al.~(2025)}}
\newcommand{\PoneEq}[1]{\href{\urlPone}{Buscicchio, Pozzoli et al.~(2025), Eq.~(#1)}}
\newcommand{\PtwoEq}[1]{\href{\urlPtwo}{Buscicchio, Klein et al.~(2025), Eq.~(#1)}}

\newcommand{\dd}{\mathrm{d}}
\newcommand{\Pbar}{\hat{P}}
\newcommand{\Mbar}{\overline{M}}

\newtcolorbox{remark}[1][]{colback=gray!6,colframe=gray!55,fonttitle=\bfseries,title=#1,breakable}
\newtcolorbox{keybox}[1][]{colback=blue!4,colframe=blue!45!black,fonttitle=\bfseries,title=#1,breakable}

\title{Chunked inference of the LISA Galactic foreground:\\ a step-by-step guide}
\author{Shri Krishna Anand}
\date{\today}

\begin{document}
\maketitle

\begin{abstract}
\noindent
This note walks through the time--frequency (``chunked'') analysis of the Galactic
foreground (GF) presented in \Pone, and through the envelope model of \Ptwo\ that it
relies on. We follow the pipeline in the order in which it is executed: frequency
grid, stationary spectra, orbital modulation, simulated data, coarse graining,
Gamma likelihood and sampling. Every equation taken from the two papers is reproduced
here, with a link to the paper and its original equation number, so the note can
be read on its own. Where our reading differs from the printed equations, we say so
explicitly.
\end{abstract}

\tableofcontents

%============================================================================
\section{The idea}
%============================================================================

The Galactic foreground is the incoherent superposition of millions of unresolved
double white-dwarf binaries in the Milky Way. Its power spectrum is not constant in
time. The Galaxy is concentrated towards the Galactic centre, and LISA's antenna
pattern sweeps across it as the constellation orbits the Sun and rotates. The power
in each TDI channel therefore rises and falls with a one-year period, twice per year.
A process whose statistics vary periodically in time is called
\emph{cyclostationary}.

The modulation is \emph{slow}: it changes appreciably over weeks to months. This
suggests the following strategy:
\begin{enumerate}
  \item split the data stream into $N_c$ non-overlapping \emph{chunks} of duration
        $T$ (one to a few weeks);
  \item inside each chunk, treat the data as \emph{stationary}, with a PSD whose
        galactic part is scaled by the modulation averaged over that chunk;
  \item treat different chunks as independent, so the total log-likelihood is a sum
        over chunks.
\end{enumerate}
The same construction handles \emph{data gaps}. Chunks are simply placed in the
stretches of data between gaps, and no special treatment of the gaps is needed.

The total log-likelihood is (\PoneEq{1})
\begin{equation}
  \log\mathcal{L}(\boldsymbol{d}\mid\theta)
  = \sum_{c=1}^{N_c} \log\mathcal{L}^{\rm chunk}(\boldsymbol{d}^c\mid\theta),
  \label{eq:total}
\end{equation}
where $\boldsymbol{d}^c$ is the data of chunk $c$ and $\theta$ the full set of model
parameters. The rest of this note explains what goes into a single
$\log\mathcal{L}^{\rm chunk}$.

\paragraph{Notation.}
\begin{center}
\begin{tabular}{@{}ll@{}}
\toprule
$c = 1,\dots,N_c$ & chunk index \\
$j\in\{A,E\}$ & TDI channel \\
$k = 1,\dots,n_f$ & fine (full-resolution) frequency bin \\
$m = 1,\dots,n_g$ & coarse (log-spaced) frequency bin \\
$l = 1,\dots,N$ & time sample within a chunk \\
$T,\ \dd t$ & chunk duration, sampling cadence \\
$N = T/\dd t,\ f_s = 1/\dd t$ & samples per chunk, sampling frequency \\
$t^c_S,\ t^c_E$ & start and end time of chunk $c$ \\
$S^c_{jk}(\theta)$ & model PSD of channel $j$, chunk $c$, frequency $f_k$ \\
\bottomrule
\end{tabular}
\end{center}

\paragraph{Overview.} \Cref{fig:flow} shows the whole pipeline. There are two
branches. The \emph{data} branch (left) is executed \emph{once}, with the injected
parameters. The \emph{model} branch (right) is executed at \emph{every} step of the
sampler, with the trial parameters $\theta$. The two meet in the likelihood.

\begin{figure}[t]
\centering
\begin{tikzpicture}[
  node distance=4mm and 8mm,
  box/.style={draw,rounded corners=2pt,align=center,font=\small,text width=6.2cm,inner sep=4pt,fill=white},
  arr/.style={-{Latex[length=2mm]},thick}]
  % data branch
  \node[box] (g)  {Chunk grid: $[t^c_S,t^c_E]$, \ $f_k = k/T$};
  \node[box,below=of g] (s)  {Stationary spectra $S^{\rm GF}(f_k)$, $S^{(n)}_j(f_k)$\\ (\S\ref{sec:spectra})};
  \node[box,below=of s] (m)  {Chunk-averaged modulation $\Mbar^c_j$\\ (\S\ref{sec:modulation})};
  \node[box,below=of m] (p)  {Total PSD $S^c_{jk}$ (\S\ref{sec:psd})};
  \node[box,below=of p] (d)  {Draw $\Re\tilde d,\Im\tilde d \sim \mathcal{N}(0,S/2)$\\ (\S\ref{sec:data})};
  \node[box,below=of d] (cg) {Average $|\tilde d_{jk}|^2$ in log bins $\to \Pbar^c_{jm},\, n_{b_m}$\\ (\S\ref{sec:coarse})};
  % model branch
  \node[box,right=12mm of s] (ms) {Same spectra at trial $\theta$,\\ at coarse frequencies $f_m$};
  \node[box,below=of ms] (mm) {$\Mbar^c_j(\theta)$ for every chunk};
  \node[box,below=of mm] (mp) {Model PSD $S^c_{jm}(\theta)$};
  \node[box,below=12mm of mp,fill=blue!6] (L) {Gamma likelihood $\log\mathcal{L}(\Pbar\mid\theta)$\\ (\S\ref{sec:gamma})};
  \node[box,below=of L,fill=blue!6] (smp) {Sampler (NUTS / nessai) proposes\\ the next $\theta$ (\S\ref{sec:inference})};
  % arrows
  \draw[arr] (g)--(s); \draw[arr] (s)--(m); \draw[arr] (m)--(p); \draw[arr] (p)--(d); \draw[arr] (d)--(cg);
  \draw[arr] (ms)--(mm); \draw[arr] (mm)--(mp); \draw[arr] (mp)--(L);
  \draw[arr] (cg.east) -- ++(4mm,0) |- (L.west);
  \draw[arr] (L)--(smp);
  \draw[arr] (smp.east) -- ++(5mm,0) |- (ms.east);
  % labels
  \node[font=\small\bfseries,above=1mm of g] {Data: once, injected $\theta$};
  \node[font=\small\bfseries,above=1mm of ms] {Model: every step, trial $\theta$};
\end{tikzpicture}
\caption{The chunked pipeline. The data are generated once and then frozen. During
sampling only the model branch is recomputed: no new random draws and no new coarse
graining.}
\label{fig:flow}
\end{figure}

%============================================================================
\section{The frequency grid of a chunk}\label{sec:grid}
%============================================================================

A chunk of duration $T$ sampled at cadence $\dd t$ contains $N = T/\dd t$ samples.
Its discrete Fourier transform lives on the grid
\begin{equation}
  f_k = \frac{k}{T}, \qquad \Delta f = \frac{1}{T},
\end{equation}
and $n_f$ denotes the number of grid points inside the analysis band
$[f_{\min},f_{\max}]$, i.e.\ $n_f \simeq (f_{\max}-f_{\min})\,T$. Longer chunks give
finer frequency resolution. With equal-duration chunks, every chunk has the
\emph{same} grid.

The Fourier amplitudes are normalised so that their squared modulus is a one-sided
PSD estimate (\PoneEq{2}):
\begin{equation}
  \tilde d^{\,c}_{jk} = \sqrt{\frac{2}{W f_s}}\ \sum_{l=1}^{N} w_l\, d^{\,c}_{jl}\,
  \exp\!\left[-\frac{2\pi\imath}{N}\,k\,l\right],
  \qquad W=\sum_{l=1}^N |w_l|^2,
  \label{eq:dft}
\end{equation}
where $w_l$ is a window function. For data generated directly in the frequency
domain (\cref{sec:data}) no window is needed and $W=N$. With this normalisation
$\langle|\tilde d^{\,c}_{jk}|^2\rangle = S^c_{jk}$.

\paragraph{Full-resolution (Whittle) likelihood.}
If the data in a chunk are Gaussian and stationary, the Fourier amplitudes at
different frequencies are independent, and the chunk likelihood is the Whittle
likelihood (\PoneEq{3}):
\begin{equation}
  \log\mathcal{L}^{\rm chunk}_W(\tilde{\boldsymbol d}^c\mid\theta) =
  -\frac12 \sum_{j\in\{A,E\}} \sum_{k=1}^{n_f}
  \left( \frac{|\tilde d^{\,c}_{jk}|^2}{S^c_{jk}(\theta)} + \log S^c_{jk}(\theta) + \log 2\pi \right).
  \label{eq:whittle}
\end{equation}
It is exact but expensive. At every sampler step it costs
$N_c\times 2\times n_f$ terms. For example, a two-week chunk has
$\Delta f\approx 0.83\,\mu$Hz, so a band of tens of mHz holds tens of thousands of
fine bins per channel and per chunk. This motivates the coarse-grained alternative
of \cref{sec:coarse}.

%============================================================================
\section{Stationary spectra: Galactic foreground shape and instrument noise}\label{sec:spectra}
%============================================================================

\paragraph{Galactic foreground spectrum.} The spectral shape of the GF is the
phenomenological model of \href{https://arxiv.org/abs/2103.14598}{Karnesis et al.~(2021)}
(\PoneEq{11}):
\begin{equation}
  S^{\rm GF}(f) = \frac{A}{2}\, f^{-7/3}\,
  \exp\!\left[-\left(\frac{f}{f_1}\right)^{\alpha}\right]
  \left[1+\tanh\!\left(\frac{f_{\rm knee}-f}{f_2}\right)\right],
  \label{eq:sgf}
\end{equation}
with amplitude $A$, slope parameter $\alpha$, and characteristic frequencies
$f_1$, $f_{\rm knee}$, $f_2$. The $f^{-7/3}$ is the inspiral spectrum of a
population of circular binaries. The $\tanh$ produces the high-frequency drop above
$f_{\rm knee}$, where binaries become sparse enough to be resolved and subtracted.

\paragraph{Instrument noise.} Two noise sources remain after TDI: test-mass (TM)
acceleration noise and optical-metrology-system (OMS) noise (\PoneEq{16} and
\PoneEq{17}):
\begin{align}
  S^n_{\rm tm}(f) &= \left(A_{\rm n}\times10^{-15}\,\mathrm{m\,s^{-2}\,Hz^{-1/2}}\right)^2
  \frac{1}{(2\pi c f)^2}
  \left[1+\left(\frac{0.4\,\mathrm{mHz}}{f}\right)^2\right]
  \left[1+\left(\frac{f}{8\,\mathrm{mHz}}\right)^4\right], \label{eq:tm}\\
  S^n_{\rm oms}(f) &= \left(P\times10^{-12}\,\mathrm{m\,Hz^{-1/2}}\right)^2
  \left(\frac{2\pi f}{c}\right)^2
  \left[1+\left(\frac{2\,\mathrm{mHz}}{f}\right)^4\right], \label{eq:oms}
\end{align}
where $c$ is the speed of light. (The paper calls the TM amplitude $A$; we write
$A_{\rm n}$ to avoid a clash with the GF amplitude.) Both components are projected
onto the A and E channels with the TDI transfer functions of
\href{https://arxiv.org/abs/2211.02539}{Nam et al.~(2023)}, giving the channel noise
PSD $S^{(n)}_{j}(f)$. The injected values are $A_{\rm n}=2.4$ and $P=7.9$.

\begin{keybox}[Why these do not depend on the chunk]
$S^{\rm GF}(f)$ and $S^{(n)}_j(f)$ depend only on frequency. They are evaluated on
the chunk grid $f_k$, which is the same for every chunk. All the time dependence of
the model comes from the modulation factor of \cref{sec:modulation}.
\end{keybox}

\begin{remark}[Remark: the GF shape evolves with observation time]
In \Pone, $f_1$ and $f_{\rm knee}$ are not constants but functions of the total
observation time (\PoneEq{12} and \PoneEq{13}):
\begin{equation}
  \log_{10} f_1 = a_1 \log_{10}(T_{\rm obs}) + b_1, \qquad
  \log_{10} f_{\rm knee} = a_{\rm knee} \log_{10}(T_{\rm obs}) + b_{\rm knee},
\end{equation}
with calibrated values $a_1=-0.25$, $b_1=-2.7$, $a_{\rm knee}=-0.27$,
$b_{\rm knee}=-2.47$. Both slopes are negative. The longer LISA observes, the more
binaries are resolved and subtracted, and the cutoff moves to lower frequency.

To mimic a global fit, where the foreground estimate is continuously updated as
individual sources are added or removed, the paper simulates the chunks
\emph{sequentially}. Each new chunk's spectrum uses the observation time
accumulated up to that chunk. The GF spectrum then drifts slowly to lower frequency
from chunk to chunk, on top of the up-and-down drift caused by the modulation. If
$T_{\rm obs}$ is held fixed, as in the rest of this note, the shape is the same for
all chunks.
\end{remark}

%============================================================================
\section{The orbital modulation}\label{sec:modulation}
%============================================================================

The modulation is derived in \Ptwo\ (main text Sec.~6 and Appendix~A) in the
\emph{low-frequency approximation}. In this limit, A and E behave as two independent
effective Michelson interferometers whose orientation follows LISA's orbit.

\subsection{Step 1: many sources at a single sky location}

Consider a large number of binaries at the same sky position: ecliptic latitude and
longitude $(b_S,l_S)$. Their inclinations $\iota$ and polarisation angles $\psi$ are
random. The squared response is averaged over the gravitational-wave period and over
these angles (\PtwoEq{27}):
\begin{equation}
  \langle h_{A,E}\rangle^2(t,b_S,l_S) = \frac{1}{T}\int_0^T \dd t'
  \int_{-1}^{1}\dd\cos\iota \int_0^{2\pi}\dd\psi\ h^2_{A,E}.
  \label{eq:env27}
\end{equation}
Here $T$ is the \emph{period of the gravitational wave} (minutes to hours). The
average removes the fast oscillation of the signal and keeps only its slowly
varying amplitude. The result still depends on time $t$, but only through LISA's
annual motion.

The averaged envelope is most conveniently written through the sum and difference
of the two channels. Both are short Fourier series in the orbital phase
(\PtwoEq{A.35}--\href{\urlPtwo}{(A.40)}):
\begin{align}
  h_S^2 &\equiv \langle h_A\rangle^2 + \langle h_E\rangle^2
         = \sum_{n=0}^{4} h^2_{S;(n)}(b_S)\, \cos\big(n\,\Delta l_L\big), \\
  h_D^2 &\equiv \langle h_A\rangle^2 - \langle h_E\rangle^2
         = \sum_{n=0}^{8} h^2_{D;(n)}(b_S)\, \cos\big(n\,\Delta l_L + 4\bar l_S\big),
\end{align}
with $\Delta l_L = \Omega t - l_S$, $\bar l_S = l_S + \pi/12$, $\Omega = 2\pi/{\rm yr}$.
The coefficients $h^2_{S;(n)}$, $h^2_{D;(n)}$ are polynomials in $\sin b_S$ and
$\cos b_S$ (\PtwoEq{A.41}--\href{\urlPtwo}{(A.54)}). For example,
$h^2_{S;(0)} = \tfrac{9}{320}\left(328 + 152\cos^2 b_S - 37\cos^4 b_S\right)$.
So the sum $A^2+E^2$ contains harmonics up to $4\Omega$, and the difference up to
$8\Omega$.

\subsection{Step 2: a model for the sky distribution}

The Galaxy is not a point. Its unresolved binaries are modelled with a (deliberately
simple, and admittedly unrealistic) bivariate Gaussian on the sky
(\PtwoEq{A.55}--\href{\urlPtwo}{(A.57)}):
\begin{equation}
  p(b_S,l_S) = \frac{1}{2\pi\sigma_1\sigma_2}\exp\!\left[-\frac12\,\boldsymbol\theta^T\mathbb{M}\,\boldsymbol\theta\right],
  \qquad
  \boldsymbol\theta = \begin{pmatrix} b_S-b_M\\ l_S-l_M\end{pmatrix},
\end{equation}
\begin{equation}
  \mathbb{M} =
  \begin{pmatrix}\cos\delta & \sin\delta\\ -\sin\delta & \cos\delta\end{pmatrix}
  \begin{pmatrix}\sigma_1^{-2} & 0\\ 0 & \sigma_2^{-2}\end{pmatrix}
  \begin{pmatrix}\cos\delta & -\sin\delta\\ \sin\delta & \cos\delta\end{pmatrix}.
\end{equation}
It is centred on the Galactic centre $(b_M,l_M)$, has standard deviations
$\sigma_{1,2}$ along its principal axes, and is rotated by an angle $\delta$
relative to the ecliptic.

\subsection{Step 3: integrate over the sky}

The envelope of the whole Galaxy is the single-location envelope weighted by $p$ and
integrated over the sky (\PtwoEq{A.58}):
\begin{equation}
  H^2_{S,D}(t) = \int_{\mathbb{R}}\cos b_S\,\dd b_S\int_{\mathbb{R}}\dd l_S\
  p(b_S,l_S)\, h^2_{S,D}.
\end{equation}
For simplicity the integral runs over the whole plane rather than over the sphere.
This is a good approximation as long as $p$ falls off quickly. The integral can be
done in closed form because every term of $h^2_{S,D}$ is a product of sines and
cosines of $b_S$ and $l_S$. The integral of $e^{\imath n b_S}e^{\imath m l_S}$
against a Gaussian is again a Gaussian factor (its characteristic function,
\PtwoEq{A.59}). Each harmonic is therefore damped by a factor that depends on the
widths and orientation of the Galaxy. With the shorthand (\PtwoEq{A.63}--\href{\urlPtwo}{(A.69)})
\begin{equation}
  \sigma_S^2=\sigma_1^2+\sigma_2^2,\quad
  \sigma_D^2=\sigma_1^2-\sigma_2^2,\quad
  \sigma_c^2=\sigma_D^2\cos2\delta,\quad
  \sigma_s^2=\sigma_D^2\sin2\delta,\quad
  \sigma_+^2=\sigma_S^2+\sigma_c^2,
\end{equation}
\begin{equation}
  \bar l_M = l_M + \frac{\pi}{12},\qquad \Delta\phi_L = \Omega t - l_M,
\end{equation}
the result is again a finite Fourier series in $\Delta\phi_L$
(\PtwoEq{A.61}, \href{\urlPtwo}{(A.62)}),
\begin{equation}
  H_S^2 = \frac{1}{10240}\sum_{n=0}^{4} H^2_{S;(n)},\qquad
  H_D^2 = \frac{1}{10240}\sum_{n=0}^{8} H^2_{D;(n)},
\end{equation}
where $H^2_{S;(n)}$ and $H^2_{D;(n)}$ contain the $n$-th harmonic. The first two
terms are (\PtwoEq{A.70}, \href{\urlPtwo}{(A.71)})
\begin{align}
  H^2_{S;(0)} &= e^{-\frac14\sigma_S^2-\frac14\sigma_c^2}
  \left(120636\cos b_M + 7614\,e^{-2\sigma_+^2}\cos3b_M - 666\,e^{-6\sigma_+^2}\cos5b_M\right),\\
  H^2_{S;(1)} &= \sqrt3\,e^{-\frac12\sigma_S^2}\Big[
  -31392\Big(\cosh\tfrac{\sigma_s^2}{2}\sin b_M\cos\Delta\phi_L + \sinh\tfrac{\sigma_s^2}{2}\cos b_M\sin\Delta\phi_L\Big) \nonumber\\
  &\qquad -32112\,e^{-2\sigma_+^2}\Big(\cosh\tfrac{3\sigma_s^2}{2}\sin3b_M\cos\Delta\phi_L + \sinh\tfrac{3\sigma_s^2}{2}\cos3b_M\sin\Delta\phi_L\Big) \nonumber\\
  &\qquad -720\,e^{-6\sigma_+^2}\Big(\cosh\tfrac{5\sigma_s^2}{2}\sin5b_M\cos\Delta\phi_L + \sinh\tfrac{5\sigma_s^2}{2}\cos5b_M\sin\Delta\phi_L\Big)\Big].
\end{align}
The remaining terms (\PtwoEq{A.72}--\href{\urlPtwo}{(A.83)}) follow the same pattern:
a damping exponential, times coefficients in $b_M$, times $\cos$ and $\sin$ of
$n\Delta\phi_L$ (for $H_S$) or of $n\Delta\phi_L+4\bar l_M$ (for $H_D$). Finally, the
envelopes of the two channels are (\PtwoEq{A.84}, \href{\urlPtwo}{(A.85)})
\begin{equation}
  \boxed{\ h^2_{A,0}(t) = \tfrac12\left(H_S^2 + H_D^2\right),\qquad
          h^2_{E,0}(t) = \tfrac12\left(H_S^2 - H_D^2\right).\ }
  \label{eq:AE}
\end{equation}
These are the squared modulation functions $M_A^2(t)$ and $M_E^2(t)$ used in \Pone.
They describe the \emph{power} of the GF in each channel as a function of time.
\Cref{tab:params} lists the parameters under the names used by each source.

\begin{table}[h]
\centering
\begin{tabular}{@{}llll@{}}
\toprule
Meaning & Klein et al. & Pozzoli et al. & \texttt{bahamas} \\
\midrule
Galactic-centre latitude  & $b_M$ & $\sin\beta$ & \texttt{lat} \\
Galactic-centre longitude & $l_M$ & $\lambda$ & \texttt{long} \\
Variances along the axes  & $\sigma_1^2,\ \sigma_2^2$ & $\sigma_1^2,\ \sigma_2^2$ & \texttt{s1}, \texttt{s2} \\
Orientation               & $\delta$ & $\sin\psi$ & \texttt{psi} ($=\sin2\delta$) \\
\bottomrule
\end{tabular}
\caption{The five envelope parameters ($\theta_{\rm GF_{II}}$) as named in \Ptwo, in \Pone\ (which samples the latitude through its sine), and in the code.}
\label{tab:params}
\end{table}

\noindent
A few printed coefficients in App.~A of \Ptwo\ contain misprints; the code implements
the corrected expressions. Details are in
\href{https://github.com/shrikrishnaanand/bahamas_ska/blob/ska_branch/reports/PAPER_BUGS.md}{\texttt{reports/PAPER\_BUGS.md}} of the \texttt{bahamas} repository.

\subsection{Step 4: average over a chunk}

Inside a chunk the PSD is assumed stationary, so the modulation is replaced by its
average over the chunk (\PoneEq{14}):
\begin{equation}
  \Mbar^c_j = \frac{1}{t^c_E-t^c_S}\int_{t^c_S}^{t^c_E}\dd t\ M_j^2(t\mid\theta_{\rm GF_{II}}).
  \label{eq:mbar}
\end{equation}
This is a single number per chunk and per channel. Note that this is a \emph{second},
different average. \Cref{eq:env27} averages over the gravitational-wave period of
each binary (fast). \Cref{eq:mbar} averages the resulting slow envelope over the
duration of a chunk (weeks).

There is no need to evaluate $M_j^2(t)$ on a fine time grid and integrate it
numerically. \Cref{eq:AE} is a finite sum of terms $\cos(n\Delta\phi_L+\varphi)$ and
$\sin(n\Delta\phi_L+\varphi)$, and the average of each term over the chunk is
elementary. For $n\geq1$,
\begin{equation}
  \frac{1}{t_E-t_S}\int_{t_S}^{t_E}\dd t\,\cos(n\Omega t+\varphi)
  = \frac{\sin(n\Omega t_E+\varphi)-\sin(n\Omega t_S+\varphi)}{n\,\Omega\,(t_E-t_S)},
\end{equation}
and similarly for the sines. Replacing every harmonic by its chunk average gives
$\Mbar^c_j$ \emph{exactly} and at negligible cost. This is what \texttt{bahamas}
does.

\begin{remark}[Why the chunk average is a reasonable approximation]
The envelope is dominated by low harmonics of the one-year orbit, so it changes
little over a week or two. Treating the chunk as stationary with the mean power
$\Mbar^c_j$ discards only the variation of $M^2_j$ \emph{within} the chunk. The
longer the chunk, the larger this error. The shorter the chunk, the fewer frequency
bins it has. The chunk duration balances the two.
\end{remark}

%============================================================================
\section{The total PSD of a chunk}\label{sec:psd}
%============================================================================

Combining the pieces, the PSD of channel $j$ in chunk $c$ is (\PoneEq{10})
\begin{equation}
  S^c_{jk}(\theta) = \Mbar^c_j(\theta_{\rm GF_{II}})\, S^{\rm GF}_k(\theta_{\rm GF_I})
  + S^{(n)}_{jk}(\theta_n) + R_{jk}\, S^{\rm EF}_k(\theta_{\rm EF}).
  \label{eq:psd}
\end{equation}
The first term is the modulated Galactic foreground: a time-independent spectral
shape, rescaled chunk by chunk. The second is the stationary instrument noise. The
last is an optional stationary, isotropic extragalactic foreground, projected with a
chunk-independent response $R_{jk}$. It is only included when that component is
studied, and we set it to zero here. The parameters are grouped as
$\theta_{\rm GF_I}=\{\log_{10}A,\log_{10}f_{\rm knee},\log_{10}f_1,\log_{10}f_2\}$
(spectral shape), $\theta_{\rm GF_{II}}$ (envelope, \cref{tab:params}) and
$\theta_n=\{A_{\rm n},P\}$ (noise).

%============================================================================
\section{Simulating the data}\label{sec:data}
%============================================================================

Both the GF and the instrument noise are modelled as zero-mean Gaussian processes.
The GF is the sum of very many independent sources, so the central limit theorem
applies. The PSD fixes the \emph{variance} at each frequency, not the values of the
data. A data set is one random realisation consistent with that variance.

\Pone\ generates the data directly in the frequency domain. For each chunk $c$,
channel $j$ and fine frequency $f_k$, the real and imaginary parts of the Fourier
amplitude are drawn independently:
\begin{equation}
  \Re\,\tilde d^{\,c}_{jk},\ \Im\,\tilde d^{\,c}_{jk} \ \sim\ \mathcal{N}\!\left(0,\ \tfrac12 S^c_{jk}(\theta_{\rm inj})\right),
  \label{eq:draw}
\end{equation}
where $\theta_{\rm inj}$ are the injected parameters. $\tilde d^{\,c}_{jk}$ is a
\emph{circular} complex Gaussian with zero mean, and its squared modulus has
expectation
\begin{equation}
  \big\langle |\tilde d^{\,c}_{jk}|^2\big\rangle
  = \big\langle (\Re\tilde d)^2\big\rangle + \big\langle (\Im\tilde d)^2\big\rangle
  = S^c_{jk}.
\end{equation}
The amplitudes themselves average to zero and carry no information. The
information sits in their \emph{squared moduli}, whose mean is the PSD. This is
why everything that follows works with $|\tilde d|^2$.

Since each chunk is drawn independently, the chunks are independent by
construction, and no window is needed ($W=N$ in \cref{eq:dft}). The data are
generated once and then kept fixed for the whole inference.

%============================================================================
\section{Coarse graining: averaged periodograms}\label{sec:coarse}
%============================================================================

To reduce the cost of \cref{eq:whittle}, the $n_f$ fine frequencies of each chunk
are grouped into $n_g\ll n_f$ \emph{logarithmically spaced} intervals. The interval
edges are (\PoneEq{5}--\href{\urlPone}{(8)})
\begin{align}
  \Delta &= \frac{\log_{10}f_{\max} - \log_{10}f_{\min}}{n_g},\\
  \log_{10} f_m &= \log_{10} f_{\min} + m\,\Delta,\\
  \mathcal{F}_m &= \{\,k \mid f_m < f_k < f_{m+1}\,\},\\
  n_{b_m} &= |\mathcal{F}_m|.
\end{align}
Here $\mathcal{F}_m$ is the set of fine frequencies that fall in the $m$-th interval,
and $n_{b_m}$ is how many there are.\footnote{\Pone\ prints the third line as
$\mathcal{F}_k=\{m\mid f_m<f_k<f_{m+1}\}$. The roles of $k$ and $m$ are swapped
there; the version above is the intended one.} Within each interval, the squared
amplitudes are replaced by their mean, the \emph{averaged periodogram}
(\PoneEq{4}):
\begin{equation}
  \Pbar^{\,c}_{jm} = \frac{1}{n_{b_m}}\sum_{k\in\mathcal{F}_m} |\tilde d^{\,c}_{jk}|^2 .
  \label{eq:phat}
\end{equation}
After this step each chunk and channel is described by $n_g$ numbers
$\Pbar^{\,c}_{jm}$ instead of $n_f$ complex amplitudes. Each $\Pbar^{\,c}_{jm}$ is an
estimate of the PSD in its interval, whose scatter shrinks as $n_{b_m}$ grows.

Three practical points:
\begin{itemize}
  \item Because the intervals have equal \emph{logarithmic} width, they are narrow at
        low frequency and wide at high frequency. $n_{b_m}$ is therefore not
        $n_f/n_g$ but grows roughly in proportion to $f$: a few fine bins near
        $f_{\min}$, hundreds at the high end. Intervals that contain no fine
        frequency are discarded.
  \item Averaging assumes that the PSD is nearly constant across an interval. With
        $n_g=1000$ (\Pone) each interval is only a fraction of a percent to about a
        percent wide in frequency, so this holds wherever the spectrum is smooth.
  \item The model PSD is evaluated at one representative frequency per interval
        (in \texttt{bahamas}, the midpoint between its first and last fine
        frequency), giving $S^c_{jm}(\theta)$.
\end{itemize}
\Pone\ uses $f_{\min}=0.1$\,mHz and $n_g=1000$.\footnote{The paper quotes
$f_{\max}=2.9$\,mHz. The \texttt{bahamas} template configuration uses
$f_{\max}=0.029$\,Hz $=29$\,mHz; it is worth confirming which value was used.}

%============================================================================
\section{The Gamma likelihood}\label{sec:gamma}
%============================================================================

\subsection{Why a Gamma distribution}

Each $|\tilde d_{jk}|^2$ is the sum of the squares of two independent zero-mean
Gaussians (real and imaginary part), each with variance $S/2$. That is
$\tfrac{S}{2}\chi^2_2$, an exponential distribution with mean $S$, which is the
special case of the Gamma distribution with shape $1$
(cf.\ \PtwoEq{21}). The Gamma family is closed under sums: independent Gamma
variables with a common scale add up to a Gamma variable whose shape is the sum of
the shapes. The averaged periodogram $\Pbar^{\,c}_{jm}$ is the mean of $n_{b_m}$
uncorrelated squared circular Gaussians, so it is Gamma distributed. (By contrast,
the amplitudes $\tilde d$ themselves are Gaussian, which is why
\cref{eq:whittle} has a Gaussian form.)

\Pone\ writes the single-chunk likelihood as (\PoneEq{9})
\begin{equation}
  \begin{split}
  \log\mathcal{L}^{\rm chunk}_G(\Pbar^{\,c}\mid\theta) =
  -\sum_{j\in\{A,E\}}\sum_{m=1}^{n_g}\Bigg[ &
  \log\Gamma\!\left(\frac{n_{b_m}}{2}\right)
  + \frac{n_{b_m}}{2}\log S^c_{jm}(\theta) \\
  &- \left(\frac{n_{b_m}}{2}-1\right)\log\Pbar^{\,c}_{jm}
  + \frac{n_{b_m}}{2}\frac{\Pbar^{\,c}_{jm}}{S^c_{jm}(\theta)}
  \Bigg].
  \end{split}
  \label{eq:gamma}
\end{equation}
The two roles are worth stating clearly:
\begin{itemize}
  \item $\Pbar^{\,c}_{jm}$ (with a hat) is the \textbf{data}: the averaged periodogram
        of \cref{eq:phat}, computed once from the simulated amplitudes.
  \item $S^c_{jm}(\theta)$ is the \textbf{model}: the PSD of \cref{eq:psd} evaluated
        at whatever $\theta$ the sampler is currently trying.
\end{itemize}
The injected (``true'') PSD does not appear in the likelihood at all. It was only
used to generate the data in \cref{sec:data}. When $\theta=\theta_{\rm inj}$, the
model coincides with it.

\subsection{Comparison with the textbook Gamma likelihood}

The Gamma density with shape $k>0$ and scale $\vartheta>0$ is
(\href{https://www.itl.nist.gov/div898/handbook/eda/section3/eda366b.htm}{NIST/SEMATECH
e-Handbook, Sec.~1.3.6.6.11}; Casella \& Berger, \emph{Statistical Inference},
Sec.~3.3)
\begin{equation}
  p(x\mid k,\vartheta) = \frac{x^{k-1}\,e^{-x/\vartheta}}{\Gamma(k)\,\vartheta^{k}},
  \qquad \langle x\rangle = k\vartheta,\qquad \mathrm{Var}(x) = k\vartheta^2 .
\end{equation}
For $n$ independent samples with the \emph{same} $k$ and $\vartheta$, the
log-likelihood is
\begin{equation}
  \ell(k,\vartheta) = (k-1)\sum_{i=1}^{n}\log x_i - \frac{1}{\vartheta}\sum_{i=1}^{n} x_i
  - n\log\Gamma(k) - n k\log\vartheta .
  \label{eq:gamma_generic}
\end{equation}
In \cref{eq:gamma}, every pair $(j,m)$ is a \emph{single} observation with its own
shape and scale. The sum over channels and bins plays the role of the sum over
samples in \cref{eq:gamma_generic}, with $n=1$ for each term. Writing one term of
\cref{eq:gamma_generic} with the overall minus sign of \cref{eq:gamma} pulled out:

\begin{center}
\renewcommand{\arraystretch}{1.4}
\begin{tabular}{@{}lll@{}}
\toprule
Textbook, one sample & Eq.~(\ref{eq:gamma}), one $(j,m)$ & Identification \\
\midrule
$x$ & $\Pbar^{\,c}_{jm}$ & the data \\
$(k-1)\log x$ & $+\left(\frac{n_{b_m}}{2}-1\right)\log\Pbar$ & shape $k = n_{b_m}/2$ \\
$-\log\Gamma(k)$ & $-\log\Gamma\!\left(\frac{n_{b_m}}{2}\right)$ & same $k$ \\
$-x/\vartheta$ & $-\frac{n_{b_m}}{2}\,\Pbar/S$ & scale $\vartheta = 2S/n_{b_m}$ \\
$-k\log\vartheta$ & $-\frac{n_{b_m}}{2}\log S$ & $=-k\log\vartheta$ up to $-\frac{n_{b_m}}{2}\log\frac{2}{n_{b_m}}$ \\
\bottomrule
\end{tabular}
\end{center}

\noindent The last entry differs only by a term that depends on $n_{b_m}$ alone. That
term does not depend on $\theta$, so \Pone\ drops it. In summary,
\begin{equation}
  \Pbar^{\,c}_{jm} \sim \mathrm{Gamma}\!\left(k = \frac{n_{b_m}}{2},\ \ \vartheta = \frac{2 S^c_{jm}(\theta)}{n_{b_m}}\right),
  \qquad \langle\Pbar\rangle = k\vartheta = S^c_{jm}(\theta).
\end{equation}

\paragraph{Interpretation.}
\begin{itemize}
  \item \emph{The model PSD is the mean, and it sets the scale.} $S$ is not the
        scale parameter itself but the mean $k\vartheta$. Since
        $\vartheta = 2S/n_{b_m}$ is proportional to $S$, changing $S$ at fixed
        $n_{b_m}$ simply stretches the distribution of $\Pbar$ without changing its
        shape.
  \item \emph{The number of averaged bins sets the shape.} The relative scatter is
        $\sqrt{\mathrm{Var}}/\langle\Pbar\rangle = 1/\sqrt{k}$. For small $n_{b_m}$
        the distribution is skewed (an exponential for $k=1$). For large $n_{b_m}$ it
        approaches a narrow Gaussian centred on $S$. High-frequency intervals, which
        average many bins, therefore constrain the model more tightly.
  \item \emph{What the sampler sees.} Only two terms depend on $\theta$:
        \begin{equation}
          \log\mathcal{L}^{\rm chunk}_G = -\sum_{j,m}\frac{n_{b_m}}{2}
          \left[\log S^c_{jm}(\theta) + \frac{\Pbar^{\,c}_{jm}}{S^c_{jm}(\theta)}\right] + \text{const}.
        \end{equation}
        Each bracket is minimised at $S=\Pbar$, and each is weighted by $n_{b_m}$.
        For $n_{b_m}=1$ this reduces to the $\theta$-dependent part of the Whittle
        likelihood, \cref{eq:whittle}. So the Gamma likelihood is the natural
        generalisation of Whittle to averaged periodograms.
\end{itemize}

\begin{remark}[Remark: counting degrees of freedom]
A complex Fourier amplitude carries \emph{two} real Gaussian degrees of freedom, so
each $|\tilde d_{jk}|^2$ is Gamma with shape $1$ (as in \PtwoEq{21}). The mean of
$n_{b_m}$ of them is therefore Gamma with shape $n_{b_m}$ and scale $S/n_{b_m}$,
with variance $S^2/n_{b_m}$. \Cref{eq:gamma} instead uses shape $n_{b_m}/2$, which
corresponds to variance $2S^2/n_{b_m}$. The $\theta$-dependent part of
\cref{eq:gamma} is then exactly one half of the exact one:
\[
  \mathcal{L}_{\rm Eq.\,(9)}(\theta) \propto \mathcal{L}_{\rm exact}(\theta)^{1/2}.
\]
This is equivalent to tempering the likelihood with inverse temperature $1/2$. The
location of the posterior is unchanged, but its widths are inflated by about
$\sqrt2$. The Whittle likelihood of \cref{eq:whittle} carries the same factor
$\tfrac12$ per complex amplitude, so the two likelihoods remain consistent with each
other. Doubling $n_{b_m}$ in \cref{eq:gamma}, and dropping the $\tfrac12$ in
\cref{eq:whittle}, gives the exact versions. We flag this as a point to discuss with
the authors.
\end{remark}

%============================================================================
\section{Inference}\label{sec:inference}
%============================================================================

Putting \cref{eq:total} and \cref{eq:gamma} together gives
$\log\mathcal{L}(\boldsymbol{\Pbar}\mid\theta)$. Bayes' theorem then gives the
posterior
$p(\theta\mid\boldsymbol{\Pbar}) \propto \mathcal{L}(\boldsymbol{\Pbar}\mid\theta)\,\pi(\theta)$.
In \Pone\ the parameter vector has 14 components: 4 for the GF spectrum
($\theta_{\rm GF_I}$), 5 for the envelope ($\theta_{\rm GF_{II}}$), 3 for the
extragalactic foreground ($\theta_{\rm EF}$), and the two noise amplitudes
$A_{\rm n}$ and $P$. All have uniform priors. Without the extragalactic foreground
there are 11.

At every step the sampler proposes a trial $\theta$, and only the model branch of
\cref{fig:flow} is recomputed: $\Mbar^c_j(\theta)$ for every chunk, the spectra at
the coarse frequencies, and the model PSD $S^c_{jm}(\theta)$. The data
$\Pbar^{\,c}_{jm}$ never change. Because every ingredient is analytic (the envelope
of \cref{eq:AE}, its chunk average and the spectra), the whole chain from $\theta$
to $\log\mathcal{L}$ is cheap and differentiable. Two kinds of sampler are used:
\begin{itemize}
  \item \textbf{Hamiltonian Monte Carlo} with the No-U-Turn Sampler (NUTS,
        \href{https://arxiv.org/abs/1111.4246}{Hoffman \& Gelman 2014}), as
        implemented in \href{https://arxiv.org/abs/1912.11554}{NumPyro}. HMC moves
        through parameter space along trajectories guided by the gradient
        $\nabla_\theta\log p(\theta\mid\boldsymbol{\Pbar})$, which is obtained by
        automatic (reverse-mode) differentiation in \texttt{jax}. No gradient has to
        be derived by hand. This makes HMC efficient even with many correlated
        parameters.
  \item \textbf{Nested sampling} with
        \href{https://arxiv.org/abs/2102.11056}{\texttt{nessai}}, which uses
        normalising flows to propose new live points. Nested sampling is not an MCMC
        method. It explores the prior volume shell by shell and returns the Bayesian
        evidence $\mathcal{Z}=\int\mathcal{L}\,\pi\,\dd\theta$ together with the
        posterior. The evidence is what allows the modulated (quasi-stationary) model
        to be compared with a stationary one, in which $\Mbar^c_j$ is constant, via
        their Bayes factor.
\end{itemize}
\Pone\ finds the same posteriors with both samplers and with both likelihoods
(Whittle and Gamma).

%============================================================================
\section{Summary}
%============================================================================

\begin{keybox}[The chunked pipeline in nine steps]
\begin{enumerate}
  \item Split the observation into $N_c$ chunks $[t^c_S,t^c_E]$ of duration $T$;
        each has the frequency grid $f_k=k/T$.
  \item Evaluate the stationary spectra $S^{\rm GF}(f_k)$ and $S^{(n)}_j(f_k)$
        (Eqs.~\ref{eq:sgf}--\ref{eq:oms}); they are the same for every chunk.
  \item For each chunk and channel, compute the chunk-averaged modulation
        $\Mbar^c_j$ (\cref{eq:mbar}) analytically from the sky-integrated envelope
        (\cref{eq:AE}).
  \item Form the total PSD $S^c_{jk}=\Mbar^c_jS^{\rm GF}_k+S^{(n)}_{jk}$ (\cref{eq:psd}).
  \item Draw $\Re\tilde d^{\,c}_{jk}$ and $\Im\tilde d^{\,c}_{jk}$ independently from
        $\mathcal{N}(0,S^c_{jk}/2)$ with the injected parameters (\cref{eq:draw}).
  \item Group the fine frequencies into $n_g$ log-spaced intervals, average
        $|\tilde d^{\,c}_{jk}|^2$ in each, and keep $\Pbar^{\,c}_{jm}$ and
        $n_{b_m}$ (\cref{eq:phat}).
  \item For a trial $\theta$, recompute $\Mbar^c_j(\theta)$ and $S^c_{jm}(\theta)$ at
        the coarse frequencies.
  \item Evaluate the Gamma likelihood (\cref{eq:gamma}) and sum it over chunks
        (\cref{eq:total}).
  \item Let the sampler (NUTS or nessai) propose the next $\theta$ and return to
        step~7.
\end{enumerate}
Steps 1--6 are done once. Steps 7--9 are repeated until the posterior is sampled.
\end{keybox}

%============================================================================
\section*{References}
%============================================================================
\begin{itemize}\sloppy
  \item R.~Buscicchio, F.~Pozzoli, D.~Chirico, A.~Sesana, \emph{The first year of LISA
        Galactic foreground},
        \href{https://arxiv.org/abs/2511.03604}{arXiv:2511.03604}.
  \item R.~Buscicchio, A.~Klein, V.~Korol, F.~Di~Renzo, C.~J.~Moore, D.~Gerosa,
        A.~Carzaniga, \emph{Test for LISA foreground Gaussianity and stationarity:
        galactic white-dwarf binaries}, Eur.\ Phys.\ J.\ C \textbf{85}, 887 (2025),
        \href{https://arxiv.org/abs/2410.08263}{arXiv:2410.08263}.
  \item N.~Karnesis, S.~Babak, M.~Pieroni, N.~Cornish, T.~Littenberg,
        \emph{Characterization of the stochastic signal originating from compact
        binaries populations as measured by LISA}, Phys.\ Rev.\ D \textbf{104},
        043019 (2021), \href{https://arxiv.org/abs/2103.14598}{arXiv:2103.14598}.
  \item G.~Casella, R.~L.~Berger, \emph{Statistical Inference}, 2nd ed., Duxbury (2002),
        Sec.~3.3.
  \item NIST/SEMATECH e-Handbook of Statistical Methods,
        \href{https://www.itl.nist.gov/div898/handbook/eda/section3/eda366b.htm}{Gamma distribution}.
  \item M.~D.~Hoffman, A.~Gelman, \emph{The No-U-Turn Sampler}, JMLR \textbf{15}, 1593
        (2014), \href{https://arxiv.org/abs/1111.4246}{arXiv:1111.4246}.
  \item D.~Phan, N.~Pradhan, M.~Jankowiak, \emph{Composable effects for flexible and
        accelerated probabilistic programming in NumPyro},
        \href{https://arxiv.org/abs/1912.11554}{arXiv:1912.11554}.
  \item M.~J.~Williams, J.~Veitch, C.~Messenger, \emph{Nested sampling with normalizing
        flows for gravitational-wave inference}, Phys.\ Rev.\ D \textbf{103}, 103006
        (2021), \href{https://arxiv.org/abs/2102.11056}{arXiv:2102.11056}.
\end{itemize}

\end{document}
